\paragraph{n8n workflow logic.}
The proposed n8n workflow implements the full experimental pipeline as a staged automation graph. Stage A initializes the experiment by defining the model configuration, frozen evidence context, output prefixes, and claim boundary. Stages B, C, and D represent the local evidence-generation phase, covering repeatability, profile-pilot, and real-variant matrix experiments, respectively. Their outputs are consolidated into a frozen local evidence base of 156 runs, which is then used as the grounding context for the AI-assisted evaluation stages. This separation ensures that model outputs are evaluated against a fixed and auditable artifact rather than against changing or implicit assumptions.

\paragraph{Role of the AI and automation components.}
Stages E and F extend the workflow with AI-assisted evaluation. Stage E provides the single-agent multi-model baseline using Claude, DeepSeek, and Mistral for evidence-summary generation, claim-boundary review, and policy-ladder checking. Stage F introduces meaningful multi-agent scenarios, including task decomposition, agent collaboration, failure recovery, verification-agent checking, and model-disagreement resolution. Stage F/F makes inter-agent communication explicit by chaining GPT-5.5 as the planner agent, Claude Fable 5 as the evidence critic, DeepSeek as the recovery agent, and Mistral as the verifier/arbiter. Each model receives the previous model output, which creates an auditable communication trace rather than an isolated single-agent decision.

\paragraph{Support for the proposed method.}
The automation pipeline supports the proposed ZKTrustLLM-Agents method by linking local evidence generation, claim-boundary control, multi-model evaluation, and multi-agent reasoning within one reproducible n8n workflow. The resulting master workflow page acts both as an executable experiment controller and as a paper-ready workflow diagram. The numerical outcomes include accuracy, valid JSON rate, latency, estimated cost, failed runs, consistency, and comparison against the single-agent baseline. The final Stage F/F execution reports 15 four-model communication records across five scenarios, with 15/15 passing records and 15/15 valid JSON chains. These results are interpreted as local evidence-grounded orchestration evidence only; they do not claim production O-RAN deployment validation, packet-capture network benchmarking, public-chain benchmarking, full media-plane QoE validation, or a 240-run six-variant ablation.

\paragraph{Figure caption.}
Full Stage A--F plus F/F n8n master orchestration workflow. The graph starts with Stage A experiment initialization and claim-boundary definition, then summarizes the Stage B--D frozen local evidence base. Stage E provides the single-agent multi-model baseline, Stage F introduces multi-agent evaluation cases, and Stage F/F shows explicit four-model inter-agent communication among GPT-5.5, Claude Fable 5, DeepSeek, and Mistral. The final aggregator produces numerical summaries and paper-ready evidence files for reporting accuracy, latency, cost, failed runs, consistency, and claim-boundary preservation.
